%(BEGIN_QUESTION)
% Copyright 2015, Tony R. Kuphaldt, released under the Creative Commons Attribution License (v 1.0)
% This means you may do almost anything with this work of mine, so long as you give me proper credit

Les veiledningsdokumentet ``LabVIEW-øvelse \#2'' som forberedelse til å gjennomføre den i datalaben, sammen med en USB-basert enhet for datainnsamling. Her skal du lære hvordan du skalerer en analog inngang, og også implementerer en sammenligningsfunksjon for å lage en alarmfunksjon i et simulert system for nivåmåling av vann.

\vskip 10pt

Merk: veiledningen du trenger for å gjøre denne øvelsen finnes i ``Instrumentation Reference'' (et sett digitale dokumentfiler) som instruktøren din har gjort tilgjengelig.

\underbar{file i01615}
%(END_QUESTION)





%(BEGIN_ANSWER)


%(END_ANSWER)





%(BEGIN_NOTES)

I denne øvelsen konfigurerer vi LabVIEW til å {\it skalere} et analogt spenningssignal slik at det representerer en væskemengde (0 til 500 gallon), og til å utløse en høy-nivåalarm med justerbart utløsningspunkt.

%INDEX% Reading assignment: LabVIEW exercise #2

%(END_NOTES)
